\clearpage
\section{Эксперименты}
\subsection{Параметры воспроизводимости}

Для обеспечения воспроизводимости экспериментов
был зафиксирован сид генерации (\texttt{seed})
на значении, равном 42.
Также во всех экспериментах с LLM
(как при типизации связей, так и при
генерации ответов)
температура генерации ($\text{temperature}$)
равнялась $0$.

% % ============================================================
\subsection{Оценка конвейера построения графа знаний}
\subsubsection{Извлечение кандидатных пар}
% % ============================================================


Извлечение может быть осуществлено тремя способами: \texttt{Strict}, \texttt{Broad} и \texttt{Combined}.
Подробное описание этих стратегий приведено в разделе 2.1.

Описание методологии оценивания изложено в разделе 1.4 «Оценка извлечения контекста».
Скрипт оценивания расположен в директории \texttt{src/eval/retrieval}.

\subsubsection*{Стратегия Strict}
В таблице~\ref{tab:strict_retrieval} приведены результаты работы стратегии
\texttt{Strict} на валидационном датасете для алгоритмов извлечения.

\begin{table}[H]
    \caption{Результат \texttt{Strict} на валидационном датасете}
    \label{tab:strict_retrieval}
    \centering
    \begin{tabular}{cccc}
        \hline
        \textbf{Precision} & \textbf{Recall} & \textbf{F\textsubscript{1}} & \textbf{Время, с} \\
        \hline
        1.0                & 0.376           & 0.547                       & 206               \\
        \hline
    \end{tabular}
\end{table}

Список извлеченных фрагментов не содержит ложноположительных результатов,
что обеспечивает идеальную точность. Однако такой подход ограничен по полноте,
так как он не способен извлечь чанки, в которых искомая сущность не упомянута явно.
Таким образом, подход задает верхнюю границу по точности и
нижнюю по полноте.
Стоит отметить, что эта стратегия извлекла все прямые упоминания сущностей.

Общее время работы стратегии на всем датасете составило 206 секунд,
что составляет около 41 секунды на подкорпус в среднем.
Кроме того, лемматизация текста заняла около 194 секунд,
а сам поиск~--- оставшиеся 12 секунд.

\subsubsection*{Стратегия Broad}

Качество работы данной стратегии зависит от значений параметров, а также используемых моделей,
поэтому для достижения наилучшего качества извлечения был проведен ряд экспериментов.

Первоначальные настройки указаны в таблице~\ref{tab:initial_hyperparameters}.
Их выбор был основан на значениях, используемых в референсных решениях
ObsidianRAG~\cite{obsidianrag} и Neo4j Graph Builder~\cite{neo4j_graph_builder},
а также на структуре валидационного датасета.
Поскольку в нем используются markdown-файлы,
имеет смысл приоритизировать разбиение текста на основе
заголовков.
Модели выбраны на основе их метрик и популярности на HuggingFace.

Инференс Embedding и Cross Encoder моделей осуществлялся локально
под управлением операционной системы Arch Linux,
оснащенной графическим ускорителем AMD Radeon RX 6700 XT с 12 ГБ видеопамяти.
Аппаратное ускорение обеспечивалось платформой ROCm (Radeon Open Compute).
Все модели использовались в формате \texttt{bfloat16},
что соответствует их оригинальной спецификации.

\begin{table}[H]
    \caption{Изначальные гиперпараметры алгоритма извлечения}
    \label{tab:initial_hyperparameters}
    \centering
    \renewcommand{\arraystretch}{1.0}
    \begin{tabularx}{\textwidth}{|>{\raggedright\arraybackslash}p{6cm}|>{\raggedright\arraybackslash}X|}
        \hline
        \textbf{Параметр}     & \textbf{Значение}                                                                                                                                                             \\
        \hline
        \multicolumn{2}{|c|}{\textbf{Эмбеддинги и разбиение текста}}                                                                                                                                          \\
        \hline
        Модель эмбеддингов    & \texttt{paraphrase-multilingual-MiniLM-L12-v2~\cite{minilm_v2}}                                                                                                               \\
        \hline
        Размер чанка          & 4000                                                                                                                                                                          \\
        \hline
        Перекрытие чанков     & 200                                                                                                                                                                           \\
        \hline
        Разделители           & \texttt{["\textbackslash{}n\#~", "\textbackslash{}n\#\#~", "\textbackslash{}n\#\#\#~", "\textbackslash{}n\textbackslash{}n", ".~", "?~", "!~", "\textbackslash{}n", "~", ""]} \\
        \hline
        \multicolumn{2}{|c|}{\textbf{Поиск}}                                                                                                                                                                  \\
        \hline
        Top-K                 & 15                                                                                                                                                                            \\
        \hline
        Вес векторного поиска & 0.5                                                                                                                                                                           \\
        \hline
        \multicolumn{2}{|c|}{\textbf{Переранжирование}}                                                                                                                                                       \\
        \hline
        Модель                & \texttt{BAAI/bge-reranker-v2-m3~\cite{bge_reranker_v2}}                                                                                                                       \\
        \hline
        Top-K                 & 5                                                                                                                                                                             \\
        \hline
        Порог                 & 0.7                                                                                                                                                                           \\
        \hline
    \end{tabularx}
\end{table}

Для поиска оптимальной конфигурации был проведен ряд экспериментов,
заключающихся в смене размера чанка, моделей эмбеддингов и переранжирования.
Их результаты представлены в
таблице~\ref{tab:broad_rest}.

Поскольку стратегия \texttt{Strict}
уже обеспечивает идеальную точность,
основной идеей экспериментов над \texttt{Broad}
было увеличение полноты при
сохранении приемлемой точности.
Допустимый порог Precision был сделан равным 0.5.
Его достижение будет первостепенной задачей,
после чего приоритет сместится на максимизацию полноты
при сохранении точности выше порога.

\begin{table}[H]
    \caption{Анализ конфигураций стратегии \texttt{Broad}.}
    \label{tab:broad_rest}
    \centering
    \begin{tabular}{ccclclccccc}
        \hline
        \textbf{\#}                         & \textbf{CS} & \textbf{CO} & \textbf{E} & \textbf{K} & \textbf{CE} & \textbf{Порог} & \textbf{P} & \textbf{R} & \textbf{F\textsubscript{1}} & \textbf{Время} \\
        \hline
        1                                   & 4000        & 200         & MiniLM     & 15         & bge-r       & 0.7            & 0.354      & 0.199      & 0.255                       & 34             \\
        2                                   & 3000        & 200         & MiniLM     & 15         & bge-r       & 0.7            & 0.395      & 0.190      & 0.256                       & 40             \\
        3                                   & 5000        & 300         & MiniLM     & 15         & bge-r       & 0.7            & 0.301      & 0.186      & 0.223                       & 27             \\
        4                                   & 3000        & 200         & MiniLM     & 15         & jina        & 0.7            & 0.553      & 0.208      & 0.302                       & 9              \\
        5                                   & 3000        & 200         & MiniLM     & 15         & jina        & 0.5            & 0.510      & 0.407      & 0.452                       & 9              \\
        6                                   & 3000        & 200         & MiniLM     & 30         & jina        & 0.5            & 0.500      & 0.453      & 0.475                       & 20             \\
        7                                   & 3000        & 200         & bge-e      & 30         & jina        & 0.5            & 0.531      & 0.468      & 0.498                       & 23             \\
        \hline
        \multicolumn{4}{l}{\textit{Strict}} &             &             &            & 1.0        & 0.376       & 0.547          & 3.4                                                                    \\
        \hline
    \end{tabular}

    \vspace{1ex}
    \raggedright\footnotesize
    \textit{Примечание:} Используемые модели: \\
    MiniLM~--- \texttt{paraphrase-multilingual-MiniLM-L12-v2}~\cite{minilm_v2}, \\
    bge-e~--- \texttt{BAAI/bge-m3}~\cite{bge_m3}, \\
    bge-r~--- \texttt{BAAI/bge-reranker-v2-m3}~\cite{bge_reranker_v2}, \\
    jina~--- \texttt{jina-reranker-v2-base-multilingual}~\cite{jina_reranker_v2}. \\
    CS~--- Chunk Size, CO~--- Chunk Overlap,
    E~--- Embedding Model, CE~--- Cross Encoder, Порог~--- минимально допустимая оценка CE,
    время в минутах
\end{table}

\subsubsection*{Влияние размера чанка}
Эксперименты 1-3 показали, что фрагмент текста не может быть слишком большим,
так как это приводит к потере деталей, что негативно сказывается на точности поиска.
Оптимальным по качеству вариантом оказался размер в 3000 символов с перекрытием в 200 символов.

Стоит отметить влияние размера чанков на время работы.
Уменьшение размера чанка ведет к увеличению числа
фрагментов, которые нужно перевести в векторное пространство,
из-за чего растет количество пар для сравнения,
и следовательно время работы.
Поэтому возникает компромисс:
увеличить размер чанков, сделав их менее детализированными, но снизив время работы,
либо уменьшить размер, увеличив время.

\subsubsection*{Роль модели переранжирования}
Модель Cross Encoder также оказалась крайне важной.
Переход с bge-r на jina увеличил $\text{F}_1$ до 0.302 (эксперимент 4),
но потребовал снижения порога с 0.7 до 0.5 для раскрытия потенциала модели (эксперимент 5).
Как итог, $\text{F}_1$ вырос до 0.452.

Изначальное значение порога оценки Cross Encoder для фильтрации результатов
было завышенным, из-за чего
часть истинно положительных кандидатов отбрасывалась.
При bge-r порог в 0.7 был оправдан, так как
модель имела низкую точность, но использование jina
решило эту проблему, повысив значение Precision с 0.395 до 0.553.

Смена модели также значительно сократила время работы алгоритма
(с 40 до 9 минут). Это объясняется вдвое меньшим количеством параметров
в модели jina (300 миллионов против 600 у bge-r),
а также более современной и вычислительно эффективной архитектурой.

Так как первостепенная цель экспериментов достигнута (порог превысил 0.5),
последующие опыты будут нацелены на максимизацию полноты
при сохранении приемлемой точности.

\subsubsection*{Количество извлекаемых чанков}
С ростом точности извлечения и значительным сокращением
времени работы появился смысл увеличить число
извлекаемых кандидатов при гибридном поиске (Top-K).
Снижение его значения не имеет особого смысла,
так как это скорее всего сильно снизит полноту,
что противоречит цели.

По результатам эксперимента 6 видно,
что увеличение Top-K спровоцировало рост шума,
а также времени работы, но увеличило полноту.

\subsubsection*{Смена модели эмбеддингов}
На замену MiniLM, переводящую тексты в 384-мерное пространство
пришла более тяжелая bge-m3, которая работает с 1024-мерными эмбеддингами.
Идея была в улучшении сохранения детализации смыслов в текстах.
Как итог, рост точности и полноты на 0.031 и 0.015 соответственно
по сравнению с результатами эксперимента 6.

\subsubsection*{Стратегия Combined}

В данной стратегии использовалась оптимальная конфигурация \texttt{Broad}.
Результаты представлены в таблице~\ref{tab:combined_retrieval}.
Видно, что \texttt{Broad} вносит весомый вклад:
полнота выросла практически вдвое (с 0.376 до 0.685),
но ожидаемо упала точность (с 1 до 0.675).
Полученные результаты подтверждают
теорию: данная стратегия действительно является компромиссом
между высокой точностью и полнотой.

\begin{table}[H]
    \caption{Результат \texttt{Combined} на валидационном датасете}
    \label{tab:combined_retrieval}
    \centering
    \begin{tabular}{clllccc}
        \hline
        \textbf{Precision} & \textbf{Recall} & \textbf{F\textsubscript{1}} \\
        \hline
        0.675              & 0.685           & 0.680                       \\
        \hline
    \end{tabular}
\end{table}

\begin{figure}[H]
    \centering
    \begin{tikzpicture}
        \begin{axis}[
                width=0.85\textwidth,
                height=6cm,
                xlabel={Номер эксперимента},
                ylabel={F\textsubscript{1}},
                xmin=0.5, xmax=7.5,
                ymin=0, ymax=0.80,
                xtick={1,2,3,4,5,6,7},
                ytick={0,0.1,0.2,0.3,0.4,0.5,0.6,0.7,0.8},
                ymajorgrids=true,
                grid style=dashed,
                legend style={at={(0.02,0.98)}, anchor=north west, font=\small},
                tick label style={font=\small},
                label style={font=\small},
            ]
            \addplot[color=blue, mark=*, thick] coordinates {
                    (1, 0.255)
                    (2, 0.256)
                    (3, 0.223)
                    (4, 0.302)
                    (5, 0.452)
                    (6, 0.475)
                    (7, 0.498)
                };
            \addlegendentry{Broad}

            \addplot[color=red, dashed, thick, mark=none] coordinates {
                    (0.5, 0.547) (7.5, 0.547)
                };
            \addlegendentry{Strict (F\textsubscript{1}=0.547)}

            \addplot[color=teal, only marks, mark=square*, mark size=4pt] coordinates {
                    (7, 0.680)
                };
            \addlegendentry{Combined (эксп. 7)}

        \end{axis}
    \end{tikzpicture}
    \caption{Динамика F\textsubscript{1} стратегии \texttt{Broad} по экспериментам.
        Красная пунктирная линия~--- уровень \texttt{Strict}-стратегии.
        Квадрат~--- результат \texttt{Combined} для лучшей конфигурации.}
    \label{fig:retrieval_f1}
\end{figure}

% % ============================================================
\subsubsection{Типизация связей}
\label{sec:llm_eval}
% % ============================================================

В рамках данной серии экспериментов
рассматривалась способность больших языковых моделей
корректно определять семантический тип связи между сущностями
на основе собранных контекстов.
Для оценки использовался вручную размеченный валидационный датасет.

Поскольку важна как точность, так и полнота классификации,
в качестве основной метрики был выбран F\textsubscript{1}-score.
Кроме того, так как в данных присутствует дисбаланс,
уместнее всего использовать Weighted F\textsubscript{1}.
Идея экспериментов заключалась в максимизации основной метрики
путем поиска оптимального промпта и модели.

Результаты опытов представлены в таблице~\ref{tab:classification_results}.
LLM модели использовались посредством API облачных платформ.
Для Qwen-3-235B~\cite{qwen_3_235b} использовался Cerebras Inference API~\cite{cerebras_api},
для Gemma-3-27B~\cite{gemma_3_27b} и Gemma-4-31B~\cite{gemma_4_31b}~--- Google Vertex AI~\cite{vertex_ai}.
Оценивание производилось согласно методологии,
описанной в разделе 1.4 «Оценка LLM-типизации».
Скрипт оценивания расположен в директории \texttt{src/eval/classification}.

\begin{table}[H]
    \caption{Weighted F\textsubscript{1} классификации типов связей}
    \label{tab:classification_results}
    \centering
    \begin{tabular}{llccccccc}
        \hline
        \textbf{Модель} & \textbf{Промпт} & \textbf{Batch}                 &
        \texttt{P}      & \texttt{H}      & \texttt{R}                     &
        \texttt{S}      & \texttt{V}      & \textbf{W. F\textsubscript{1}}                                                          \\
        \hline
        Gemma-3-27B     & base            & 1                              & 0.346 & 0.340 & 0.351 & 0.459 & 0.522 & 0.393          \\
        Gemma-3-27B     & base            & 5                              & 0.446 & 0.325 & 0.321 & 0.451 & 0.573 & 0.387          \\
        Gemma-3-27B     & base            & 10                             & 0.378 & 0.428 & 0.363 & 0.392 & 0.572 & 0.423          \\
        Qwen-3-235B     & base            & 1                              & 0.313 & 0.454 & 0.272 & 0.454 & 0.526 & 0.392          \\
        Gemma-3-27B     & few-shot        & 1                              & 0.372 & 0.436 & 0.444 & 0.341 & 0.559 & 0.445          \\
        Qwen-3-235B     & few-shot        & 1                              & 0.498 & 0.390 & 0.434 & 0.558 & 0.482 & 0.445          \\
        Gemma-4-31B     & few-shot        & 1                              & 0.856 & 0.578 & 0.627 & 0.786 & 0.758 & \textbf{0.669} \\
        \hline
    \end{tabular}

    \vspace{1ex}
    \raggedright\footnotesize
    \textit{Примечание:} Используемые промпты
    доступны в репозитории проекта по пути: \texttt{artifacts/prompts/classification.} \\
\end{table}

\subsubsection*{Влияние пакетного промптинга}
Увеличение количества пар заметок, передаваемых модели в рамках одного запроса,
в среднем снижает среднюю оценку, однако в некоторых случаях она растет.
Увеличение может быть вызвано утечкой контекста:
модель получает дополнительные данные из других элементов пакета.
Это видно в случае Personal, Hub-n-Spoke и Versus, где размер содержимого заметок мал,
и есть кластеры из тематически связных заметок.
Тем не менее, важно не перегрузить модель лишним контекстом: при классификации конкретной связи
контекст для других считается шумом и не должен быть учтен. Не все модели могут корректно
разграничить содержимое, из-за чего начинают путаться и делать ошибочные выводы, ухудшая качество классификации.
Как наиболее честный вариант был выбран размер пакета равный 1.

\subsubsection*{Влияние промпта}
Добавление в промпт примеров и описаний типов связей, которые в них используются,
увеличило значение Overall F\textsubscript{1} как для Gemma-3-27B с 0.392 до 0.423 ~(+0.031),
так и для Qwen-3-235B с 0.392 до 0.445~(+0.135).
Это подтверждает чувствительность моделей
к инструкциям и примерам: они помогают им лучше понять специфику задачи
и не уйти в сторону при рассуждении.
Важно точно определить каждый тип связи для устранения неоднозначностей.
Кроме того, результаты показывают,
что качество более крупной модели больше зависит от промпта.

\subsubsection*{Влияние выбора модели}
Эксперименты показали, что
эффективность архитектуры важнее числа параметров.
Новая Gemma-4-31B (31 млрд параметров) значительно превосходит
по качеству Qwen-3-235B (235 млрд параметров),
а Gemma-3-27B (27 млрд параметров) оказалась в ней наравне.

\subsubsection*{Анализ ошибок классификации}
Несмотря на рост общего значения Weighted F\textsubscript{1},
имеется несколько общих ошибок:
\begin{itemize}
    \item Злоупотребление универсальными типами.
          Тип \texttt{mentions} является самым популярным в датасете
          (40 связей, F\textsubscript{1} = 0.23--0.39).
          Он играет роль резервного класса:
          он назначается в случае,
          когда связь есть, но невозможно однозначно назначить другой класс.
          Модель может слишком часто назначать этот тип из-за его универсальности.
          Эта проблема решается дополнительными инструкциями и примерами в промпте,
          явно указывающими смысл этой связи.
    \item Использование близких по смыслу связей.
          Классы \texttt{encompasses}, \texttt{contains}, \texttt{incorporates} и другие
          имеют высокое семантическое сходство.
          Большое их количество создает путаницу:
          модель не способна выявить семантической разницы между ними,
          поэтому с высокой вероятностью может использовать неверный тип.
          Проблема обостряется в реальных условиях неточного контекста,
          который создает алгоритм извлечения контекста.
          Эффективным решением является объединение
          подобных классов в один более универсальный.
\end{itemize}

\subsection{Сравнение сгенерированных ссылок с эталонными}

Перед оценкой GraphRAG на сгенерированных ссылках необходимо их проанализировать.
Оценка проводилась в двух режимах: без учета типа (оценивается только факт наличия связи между заметками)
и с учетом типа (оценивается как наличие связи, так и правильность предсказанного типа)
для анализа качества ссылок для использования в каждом типе обхода.
Использовались лучшие конфигурации извлечения (стратегия \texttt{Combined}) и типизации
(Gemma-4-31B, few-shot промпт, \texttt{batch=1}).

\begin{table}[H]
    \caption{Оценка качества сгенерированных связей: факт наличия (Untyped) и точное совпадение типа (Typed)}
    \label{tab:generated_links_eval}
    \centering
    \begin{tabular}{lcccccc}
        \hline
        \multirow{2}{*}{\textbf{Подкорпус}} & \multicolumn{3}{c}{\textbf{Без учета типа (Untyped)}} & \multicolumn{3}{c}{\textbf{С учетом типа (Typed)}}                                                                                       \\
        \cline{2-7}
                                            & \textbf{P}                                            & \textbf{R}                                         & \textbf{F\textsubscript{1}} & \textbf{P} & \textbf{R} & \textbf{F\textsubscript{1}} \\
        \hline
        \texttt{hub\_n\_spoke}              & 0.810                                                 & 0.883                                              & 0.845                       & 0.643      & 0.701      & 0.671                       \\
        \texttt{personal}                   & 0.909                                                 & 0.435                                              & 0.588                       & 0.909      & 0.435      & 0.588                       \\
        \texttt{realistic}                  & 0.771                                                 & 0.810                                              & 0.790                       & 0.687      & 0.722      & 0.704                       \\
        \texttt{sequence}                   & 0.733                                                 & 0.957                                              & 0.830                       & 0.533      & 0.696      & 0.604                       \\
        \texttt{versus}                     & 0.727                                                 & 0.909                                              & 0.808                       & 0.545      & 0.682      & 0.606                       \\
        \hline
        \textbf{Взвеш. усред.}              & 0.776                                                 & 0.829                                              & 0.802                       & 0.635      & 0.679      & 0.656                       \\
        \hline
    \end{tabular}
\end{table}

Точность извлеченных ссылок без учета их типа находится на высоком уровне:
$\text{F}_1 = 0.802$, при этом полнота также имеет близкое к идеальному значение
$\text{R} = 0.829$, что говорит о способности алгоритма находить
практически все ссылки.
Наилучшие результаты наблюдаются в подкорпусах
\texttt{sequence} ($\text{R} = 0.957$) и \texttt{versus} ($\text{R} = 0.909$).

Введение типов ссылок как критерия оценивания
снизило значения метрик: $\text{F}_1$ упал до $0.656$.
Исключением является образец \texttt{personal}, в котором метрики не изменились.
Это означает, что алгоритм безошибочно классифицировал все связи.
Такой результат был ожидаем, он напрямую коррелирует с результатами раздела~\ref{sec:llm_eval}.

% % ============================================================
\subsection{Оценка конвейера GraphRAG}
\label{sec:graphrag_eval}
% % ============================================================

В таблице~\ref{tab:graphrag_default_params} представлены первоначальные параметры алгоритмов.
Они выбраны на основе количества валидационных данных и референсного решения.
Значения подобраны так, чтобы заставить алгоритмы обходить граф,
а не собрать контекст всех заметок разом.

\begin{table}[H]
    \caption{Значения гиперпараметров алгоритма GraphRAG по умолчанию}
    \label{tab:graphrag_default_params}
    \centering
    \renewcommand{\arraystretch}{1.0}
    \begin{tabularx}{\textwidth}{|>{\raggedright\arraybackslash}X|l|c|}
        \hline
        \textbf{Описание параметра}            & \textbf{Переменная}       & \textbf{Значение} \\
        \hline
        Максимальная глубина обхода графа      & \texttt{MAX\_HOPS}        & 2                 \\
        \hline
        Порог отсечения узлов по релевантности & \texttt{SCORE\_THRESHOLD} & 0.2               \\
        \hline
        Количество первоначальных заметок      & \texttt{TOP\_K\_SEED}     & 3                 \\
        \hline
        Итоговое количество чанков в контексте & \texttt{TOP\_K\_CONTEXT}  & 10                \\
        \hline
    \end{tabularx}
\end{table}

Оценка проводилась в соответствии с методологией,
описанной в разделе 1.4 «Полная оценка GraphRAG».
Скрипты оценивания расположены в директории \texttt{src/eval/ragas}.

\subsubsection{Анализ референсных решений}
\label{sec:reference_results}

Для оценки разработанных алгоритмов было проведено
их комплексное сравнение с ObsidianRAG и Neo4j
(Graph Builder и его же интегрированная реализация RAG).
В качестве модели для генерации ответов была выбрана Gemma-3-27B,
для оценки Gemma-4-31B.

\begin{table}[H]
    \caption{Результаты решения ObsidianRAG}
    \label{tab:reference_results_obsidianrag}
    \centering
    \begin{tabular}{lcccc}
        \hline
        \textbf{Подкорпус}     & \textbf{CR} & \textbf{CP} & \textbf{F} & \textbf{AC} \\
        \hline
        \texttt{hub\_n\_spoke} & 0.802       & 0.701       & 0.955      & 0.542       \\
        \texttt{personal}      & 0.625       & 0.615       & 1.000      & 0.424       \\
        \texttt{realistic}     & 0.886       & 0.707       & 0.898      & 0.505       \\
        \texttt{sequence}      & 1.000       & 0.824       & 0.967      & 0.582       \\
        \texttt{versus}        & 0.617       & 0.467       & 0.883      & 0.428       \\
        \hline
        \textbf{Среднее}       & 0.786       & 0.663       & 0.941      & 0.496       \\
        \hline
    \end{tabular}
\end{table}

Результаты оценки референсного решения ObsidianRAG приведены в таблице~\ref{tab:reference_results_obsidianrag}.
Явно видны недостатки логики расширения контекста: в
\texttt{personal} и \texttt{versus} наблюдается низкий CR (0.625 и 0.617 соответственно).
В этих образцах есть два кластера, которые практически друг с другом не связаны.
Наивный сбор контекста соседних узлов не позволит получить информацию из другого кластера,
а лишь увеличит шум, что видно по низким значениям CP
(0.467 в случае versus и 0.615 в случае personal).
В других образцах, где заметки связаны большим числом ссылок, CR не ниже 0.8, а CP не ниже 0.7,
что говорит о хорошей работе алгоритма.

Как результат, следует и низкий AC:
контекст зашумлен и в нем недостаточно
данных для предоставления
полного и точного ответа.

\begin{table}[H]
    \caption{Результаты решения Neo4j на сгенерированном ГЗ}
    \label{tab:reference_results_neo4j}
    \centering
    \begin{tabular}{lcccc}
        \hline
        \textbf{Подкорпус}     & \textbf{CR} & \textbf{CP} & \textbf{F} & \textbf{AC} \\
        \hline
        \texttt{hub\_n\_spoke} & 1.000       & 0.782       & 0.951      & 0.598       \\
        \texttt{personal}      & 1.000       & 0.875       & 0.984      & 0.612       \\
        \texttt{realistic}     & 0.985       & 0.842       & 0.975      & 0.620       \\
        \texttt{sequence}      & 0.980       & 0.855       & 0.971      & 0.640       \\
        \texttt{versus}        & 0.920       & 0.751       & 0.905      & 0.540       \\
        \hline
        \textbf{Среднее}       & 0.977       & 0.821       & 0.957      & 0.602       \\
        \hline
    \end{tabular}
\end{table}

В таблице~\ref{tab:reference_results_neo4j} представлены результаты Neo4j,
полученные на сгенерированном Neo4j Graph Builder ГЗ.

Решение имеет значительно более высокие значения CR,
что обусловлено более продвинутой логикой сбора контекста.
Среднее значение CR достигло $0.977$,
в то время как у ObsidianRAG оно лишь $0.786$.
Также виден большой прирост метрики на подкорпусах
\texttt{personal} (с $0.625$ до $1.000$) и
\texttt{versus} (с $0.617$ до $0.920$).

Кроме того, CP также вырос~---
в среднем $0.821$ против $0.663$ у ObsidianRAG.
Это говорит об отличной способности Neo4j фильтровать контекст.

Как следствие, увеличились значения AC:
среднее значение выросло на $0.106$ балла (с $0.496$ до $0.602$).
Прирост особенно заметен на проблемных для ObsidianRAG подкорпусах
\texttt{personal} ($+0.188$) и \texttt{versus} ($+0.112$).

Во всех случаях наблюдается высокий Answer Faithfulness:
модель не пытается додумать ответ, а использует приведенные данные.

\subsubsection{Сравнение собственных алгоритмов с аналогами}

Поскольку основным преимуществом алгоритмов, представленных в данной работе,
является многошаговый обход, основная цель которого увеличение CR,
имеет смысл сфокусироваться только на полноте контекста и качестве ответа
при сравнении разработанных алгоритмов с эталонным.
Сравнение представлено в таблице~\ref{tab:complex_graphrag_comparison}.

\begin{table}[H]
    \caption{Сравнение собственных алгоритмов с аналогами на размеченном (GT) и сгенерированном (GEN) графах.}
    \label{tab:complex_graphrag_comparison}
    \centering
    \renewcommand{\arraystretch}{1.1}
    \setlength{\tabcolsep}{4pt}
    \begin{tabular}{l|l|cc|cc|cc|cc}
        \hline
        \multirow{2}{*}{\textbf{Подкорпус}} & \multirow{2}{*}{\textbf{Данные}} & \multicolumn{2}{c|}{\textbf{Ref}} & \multicolumn{2}{c|}{\textbf{Base}} & \multicolumn{2}{c|}{\textbf{Typed}} & \multicolumn{2}{c}{\textbf{Neo4j}}                                                         \\
        \cline{3-10}
                                            &                                  & \textbf{CR}                       & \textbf{AC}                        & \textbf{CR}                         & \textbf{AC}                        & \textbf{CR} & \textbf{AC} & \textbf{CR} & \textbf{AC} \\
        \hline
        \multirow{2}{*}{\texttt{hub}}
                                            & GT                               & 0.802                             & 0.542                              & 1.000                               & 0.587                              & 0.958       & 0.565       & ---         & ---         \\
                                            & GEN                              & 0.750                             & 0.516                              & 0.958                               & 0.535                              & 0.958       & 0.561       & 1.000       & 0.598       \\
        \hline
        \multirow{2}{*}{\texttt{personal}}
                                            & GT                               & 0.625                             & 0.424                              & 1.000                               & 0.597                              & 1.000       & 0.579       & ---         & ---         \\
                                            & GEN                              & 0.625                             & 0.429                              & 1.000                               & 0.602                              & 1.000       & 0.642       & 1.000       & 0.612       \\
        \hline
        \multirow{2}{*}{\texttt{realistic}}
                                            & GT                               & 0.886                             & 0.505                              & 0.977                               & 0.606                              & 0.977       & 0.601       & ---         & ---         \\
                                            & GEN                              & 0.864                             & 0.522                              & 0.977                               & 0.637                              & 0.977       & 0.602       & 0.985       & 0.620       \\
        \hline
        \multirow{2}{*}{\texttt{sequence}}
                                            & GT                               & 1.000                             & 0.582                              & 0.925                               & 0.580                              & 0.975       & 0.631       & ---         & ---         \\
                                            & GEN                              & 0.975                             & 0.555                              & 0.925                               & 0.614                              & 0.925       & 0.600       & 0.980       & 0.640       \\
        \hline
        \multirow{2}{*}{\texttt{versus}}
                                            & GT                               & 0.617                             & 0.428                              & 0.900                               & 0.524                              & 0.788       & 0.507       & ---         & ---         \\
                                            & GEN                              & 0.608                             & 0.388                              & 0.900                               & 0.550                              & 0.788       & 0.520       & 0.920       & 0.540       \\
        \hline
        \hline
        \multirow{2}{*}{\textbf{Среднее}}
                                            & GT                               & 0.786                             & 0.496                              & 0.960                               & 0.579                              & 0.940       & 0.577       & ---         & ---         \\
                                            & GEN                              & 0.764                             & 0.482                              & 0.952                               & 0.587                              & 0.930       & 0.585       & 0.977       & 0.602       \\
        \hline
    \end{tabular}
\end{table}

\subsubsection*{Оценка GraphRAG на размеченных данных}

Видно, что модифицированная логика обхода действительно
увеличила количество собираемой информации по сравнению с ObsidianRAG,
позволив получить практически идеальную полноту контекста.
Это оказалось особенно заметно в случае \texttt{personal} и \texttt{versus},
где изначально значение CR было самым низким.
В их случае прирост составил 0.375 и 0.283 соответственно в случае Base,
0.375 и 0.171 в случае Typed.

Дополнительный контекст помог языковой модели дать более точный ответ,
что также увеличило средние значение AC с
$0.496$ до $0.579$ (Base) и $0.577$ (Typed)

% % ============================================================
\subsubsection*{Оценка GraphRAG на сгенерированном графе}

Для получения первоначальных данных и оценки качества
работы алгоритмов на автоматически построенных ГЗ
в следующем эксперименте размеченный ГЗ был заменен
на граф, построенным пайплайном генерации связей.
Результаты приведены в таблице~\ref{tab:complex_graphrag_comparison}.

Видно, что замена графа минимально влияет на качество работы GraphRAG:
значения метрик колеблются в диапазоне сотых долей.

Переход на сгенерированные связи незначительно снизил CR:
для базового алгоритма среднее значение упало
с $0.960$ до $0.952$, а для типизированного~— с $0.940$ до $0.930$.
Это падение можно
объяснить пропуском неочевидных связей конвейером построения ГЗ,
из-за чего усложняется навигация по графу.

Интересным наблюдением является рост AC
(с $0.579$ до $0.587$ для Base и с $0.577$ до $0.585$ для Typed).
Особенно заметно увеличение в
подкорпусе \texttt{personal} для Typed-обхода (с $0.579$ до $0.642$)
и в \texttt{realistic} для Base-обхода (с $0.606$ до $0.637$).
Возможно, что сгенерированные связи в некоторых случаях
более эффективны для сбора релевантного контекста,
так как пайплайн мог иначе типизировать связи или создать новые,
которые помогли извлечь нужные данные.

Алгоритм Base показал себя лучше в большинстве случаев
как при сборе контекста, так и при генерации ответов.

Результаты работы предложенных алгоритмов оказались сопоставимы
с Neo4j: в среднем алгоритм Base уступает ему всего на $0.025$ по полноте и на $0.015$
по точности ответа, алгоритм Typed на $0.047$ и $0.017$  соответственно.
Это доказывает эффективность разработанных алгоритмов.

% ============================================================
\subsubsection{Влияние промпта и языковой модели}
\label{sec:llm_prompt_effect}

Для исследования влияния различных факторов и компонент GraphRAG на качество ответа LLM
была проведена серия экспериментов над решениями, представленными в данной работе.
Все замеры проводились на размеченных данных.

\subsubsection*{Базовый промпт против few-shot}
На основе экспериментов над LLM-классификацией связей был сделан вывод о важности
указания точных инструкций и конкретных примеров в промпте.
В базовом промпте уже есть основные указания,
такие как придерживание предоставленного контекста и ответ на языке запроса.
К ним были добавлены конкретные развернутые примеры и
требования модели к рассуждениям перед ответом.
Модели для оценки и генерации остались теми же.

\begin{table}[H]
    \caption{Сравнение базового и few-shot промптов на алгоритмах Base и Typed (метрика AC, модель Gemma-3-27B)}
    \label{tab:prompt_effect_full_versus}
    \centering
    \begin{tabular}{llcccccc}
        \hline
        \textbf{Алгоритм}      & \textbf{Промпт} & \texttt{H} & \texttt{P} & \texttt{R} & \texttt{S} & \texttt{V} & \textbf{AC avg} \\
        \hline
        \multirow{2}{*}{Base}  & Базовый         & 0.587      & 0.597      & 0.606      & 0.580      & 0.524      & 0.579           \\
                               & Few-shot        & 0.595      & 0.572      & 0.604      & 0.618      & 0.569      & 0.592           \\
        \hline
        \multirow{2}{*}{Typed} & Базовый         & 0.565      & 0.579      & 0.601      & 0.631      & 0.507      & 0.577           \\
                               & Few-shot        & 0.600      & 0.611      & 0.606      & 0.629      & 0.555      & 0.580           \\
        \hline
    \end{tabular}

    \vspace{1ex}
    \raggedright\footnotesize
    \textit{Примечание:} Используемые промпты
    доступны в репозитории проекта по пути: \texttt{artifacts/prompts/graphrag.} \\
\end{table}

Новый промпт продемонстрировал прирост AC в среднем лишь +0.13.
Прирост оказался ниже, чем в предыдущих сравнениях, вероятно,
из-за корректно сформулированных инструкций в базовом промпте.
Можно сделать вывод,
что наличие примеров все же дает небольшой прирост корректности ответа языковой модели.

\subsubsection*{Анализ влияния языковой модели}

В этом эксперименте
модель для генерации будет заменена на Gemma-4-31B, а
моделью-судьей будет Gemini-3-flash~\cite{gemini_3_flash}.

\begin{table}[H]
    \caption{Сравнение моделей Gemma-3-27B и Gemma-4-31B (few-shot промпт)}
    \label{tab:model_comparison_2_3}
    \centering
    \begin{tabular}{llccc|ccc}
        \hline
        \multirow{2}{*}{\textbf{Датасет}}   & \multirow{2}{*}{\textbf{Алгоритм}} & \multicolumn{3}{c|}{\textbf{Эксп. 2 (Gemma-3-27B)}} & \multicolumn{3}{c}{\textbf{Эксп. 3 (Gemma-4-31B)}}                                                                            \\
        \cline{3-8}
                                            &                                    & \textbf{CP}                                         & \textbf{CR}                                        & \textbf{AC} & \textbf{CP}        & \textbf{CR}        & \textbf{AC}      \\
        \hline
        \multirow{2}{*}{\texttt{hub}}       & Base                               & 0.749                                               & 1.000                                              & 0.595       & 0.660 $\downarrow$ & 0.875 $\downarrow$ & 0.698 $\uparrow$ \\
                                            & Typed                              & 0.796                                               & 0.958                                              & 0.600       & 0.707 $\downarrow$ & 0.833 $\downarrow$ & 0.709 $\uparrow$ \\
        \hline
        \multirow{2}{*}{\texttt{personal}}  & Base                               & 0.893                                               & 1.000                                              & 0.572       & 0.750 $\downarrow$ & 1.000 $=$          & 0.796 $\uparrow$ \\
                                            & Typed                              & 0.892                                               & 1.000                                              & 0.611       & 0.729 $\downarrow$ & 0.917 $\downarrow$ & 0.770 $\uparrow$ \\
        \hline
        \multirow{2}{*}{\texttt{realistic}} & Base                               & 0.848                                               & 0.977                                              & 0.604       & 0.426 $\downarrow$ & 1.000 $\uparrow$   & 0.738 $\uparrow$ \\
                                            & Typed                              & 0.851                                               & 0.977                                              & 0.606       & 0.491 $\downarrow$ & 0.955 $\downarrow$ & 0.792 $\uparrow$ \\
        \hline
        \multirow{2}{*}{\texttt{sequence}}  & Base                               & 0.879                                               & 0.925                                              & 0.618       & 0.434 $\downarrow$ & 0.765 $\downarrow$ & 0.734 $\uparrow$ \\
                                            & Typed                              & 0.878                                               & 0.975                                              & 0.629       & 0.429 $\downarrow$ & 0.765 $\downarrow$ & 0.706 $\uparrow$ \\
        \hline
        \multirow{2}{*}{\texttt{versus}}    & Base                               & 0.734                                               & 0.900                                              & 0.569       & 0.519 $\downarrow$ & 0.655 $\downarrow$ & 0.623 $\uparrow$ \\
                                            & Typed                              & 0.696                                               & 0.813                                              & 0.555       & 0.547 $\downarrow$ & 0.568 $\downarrow$ & 0.576 $\uparrow$ \\
        \hline
        \multirow{2}{*}{\textbf{Среднее}}   & Base                               & 0.821                                               & 0.960                                              & 0.592       & 0.558 $\downarrow$ & 0.859 $\downarrow$ & 0.718 $\uparrow$ \\
                                            & Typed                              & 0.823                                               & 0.945                                              & 0.600       & 0.581 $\downarrow$ & 0.808 $\downarrow$ & 0.711 $\uparrow$ \\
        \hline
    \end{tabular}
\end{table}

Результаты представлены в таблице~\ref{tab:model_comparison_2_3}.
Прирост AC составил в среднем +0.128,
при этом CP упал в среднем на 0.263.
Сильнее всего CP снизился на \texttt{realistic} (с 0.848 до 0.426)
и \texttt{sequence} (с 0.879 до 0.434).
Снижение CP обусловлено более строгой оценкой модели.
Однако рост AC показывает, что более мощная модель
для генерации ответов способна абстрагироваться от шума
и дать более точный ответ.

\subsubsection*{Сравнение типизированного обхода с обычным}

У типизированного обхода выше CP (среднее: 0.581 против 0.558 у Base), но ниже
CR (0.808 против 0.859 у Base).
Эти результаты показывают влияние фильтрации заметок при обходе графа:
тематически нерелевантные файлы отсекаются, но при этом также теряется возможность
сделать дополнительный шаг через этот файл, чтобы оценить его соседей, из-за чего упускается часть информации.

В плотно связанных подкорпусах, таких как \texttt{realistic},
типизированный подход обладает большей точностью контекста, но
на разреженных графах доминирует обычный обход, видимо,
потому что он позволяет делать дополнительные шаги через нерелевантные узлы,
что иногда позволяет добраться до необходимой информации.

При этом нельзя сделать однозначный вывод о влиянии учета типов ссылок на корректность ответов:
значения никак не коррелируют.

\subsubsection{Анализ влияния графа}

Результаты в таблице~\ref{tab:exp3_vs_exp5}
уже показали устойчивость и эффективность сгенерированных ссылок.
Данный эксперимент нацелен подтвердить,
что этот результат сохраняется при использовании
более мощных и строгих LLM как для генерации ответа,
так и для оценки.

\begin{table}[H]
    \caption{Сравнение качества работы GraphRAG между размеченным графом и сгенерированным.}
    \label{tab:exp3_vs_exp5}
    \centering
    \begin{tabular}{llccc|ccc}
        \hline
        \multirow{2}{*}{\textbf{Датасет}}   & \multirow{2}{*}{\textbf{Алгоритм}} & \multicolumn{3}{c|}{\textbf{Разметка}} & \multicolumn{3}{c}{\textbf{Сгенерированные}}                                                                               \\
        \cline{3-8}
                                            &                                    & \textbf{CP}                            & \textbf{CR}                                  & \textbf{AC} & \textbf{CP}        & \textbf{CR}        & \textbf{AC}         \\
        \hline
        \multirow{2}{*}{\texttt{hub}}       & Base                               & 0.660                                  & 0.875                                        & 0.698       & 0.689 $\uparrow$   & 0.833 $\downarrow$ & 0.682 $\downarrow$  \\
                                            & Typed                              & 0.707                                  & 0.833                                        & 0.709       & 0.707 $=$          & 0.833 $=$          & 0.716 $\uparrow$    \\
        \hline
        \multirow{2}{*}{\texttt{personal}}  & Base                               & 0.750                                  & 1.000                                        & 0.796       & 0.729 $\downarrow$ & 0.958 $\downarrow$ & 0.762 $\downarrow$  \\
                                            & Typed                              & 0.729                                  & 0.917                                        & 0.770       & 0.729 $=$          & 0.917 $=$          & 0.770 $=$           \\
        \hline
        \multirow{2}{*}{\texttt{realistic}} & Base                               & 0.426                                  & 1.000                                        & 0.738       & 0.426 $=$          & 1.000 $=$          & 0.738 $=$           \\
                                            & Typed                              & 0.491                                  & 0.954                                        & 0.792       & 0.476 $\downarrow$ & 0.977 $\uparrow$   & 0.774 $\downarrow$  \\
        \hline
        \multirow{2}{*}{\texttt{sequence}}  & Base                               & 0.434                                  & 0.765                                        & 0.734       & 0.434 $=$          & 0.765 $=$          & 0.734 $=$           \\
                                            & Typed                              & 0.429                                  & 0.765                                        & 0.706       & 0.429 $=$          & 0.765 $=$          & 0.706 $=$           \\
        \hline
        \multirow{2}{*}{\texttt{versus}}    & Base                               & 0.519                                  & 0.655                                        & 0.623       & 0.519 $=$          & 0.655 $=$          & 0.623 $=$           \\
                                            & Typed                              & 0.547                                  & 0.568                                        & 0.576       & 0.547 $=$          & 0.568 $=$          & 0.576 $=$           \\
        \hline
        \multirow{2}{*}{\textbf{Среднее}}   & Base                               & 0.558                                  & 0.859                                        & 0.718       & 0.559 $\uparrow$   & 0.842 $\downarrow$ & 0.708 $\downarrow$  \\
                                            & Typed                              & 0.581                                  & 0.807                                        & 0.711       & 0.578 $\downarrow$ & 0.812 $\uparrow$   & 0.708  $\downarrow$ \\
        \hline
    \end{tabular}
\end{table}


На основе полученных результатов (таблица~\ref{tab:exp3_vs_exp5})
можно сделать вывод, что использование сгенерированных связей
практически не влияет на качество работы GraphRAG.

Падение метрики AC при переходе на автоматически построенный граф
в среднем не превышает 0.01
(с $0.718$ до $0.708$ для Base-обхода и с $0.711$ до $0.708$ для Typed-обхода).
Самое большое падение AC наблюдается в случае базового алгоритма
в образце \texttt{personal} и составляет 0.034 (с 0.796 до 0.762).

Изменения значений Context Precision и Context Recall
незначительны, что говорит о сопоставимом качестве связей
с ручной разметкой.

\subsection{Выводы и ограничения исследования}
\label{sec:conclusions}

Опыты над алгоритмом построения графа определили оптимальную конфигурацию
извлечения контекста: переход к рекурсивному алгоритму вместе с использованием
Jina как модели переранжирования позволил достичь $F_1 = 0.498$,
что в комбинации с результатами стратегии \texttt{Strict} дает
оптимальный баланс между точностью и полнотой извлеченных фрагментов.

Эксперименты над LLM-классификацией связи доказали критичность использования
few-shot промпта с четкими инструкциями, а также важность современной архитектуры языковой модели:
Gemma-4-31B показала результат Overall $F_1 = 0.669$,
превзойдя более крупную Qwen-3-235B ($F_1 = 0.445$).
Также важным оказалось объединение семантически схожих классов ссылок
для упрощения классификации языковой моделью.

Оценка предложенных алгоритмов GraphRAG через RAGAS
показала их эффективность:
по результатам видно явное превосходство над ObsidianRAG
и сопоставимое с Neo4j качество работы.
Разработанная логика сбора контекста повысила CR по всех подкорпусах вплоть до 1.0.
Сравнение типизированного обхода с обычным показало гибкость системы:
алгоритм может точно извлекать достаточный объем информации в разных сценариях.
Типизированный обход лучше всего подходит для работы с
шумными, плотно связанными графами.
Базовый, в свою очередь, демонстрирует отличные результаты на разреженных графах.

Финальное сравнение работы предложенных алгоритмов GraphRAG
на сгенерированных ссылках показало высокое качество пайплайна для построения графов знаний в Obsidian.
Конвейер устойчив к шуму и может справляться с неоднозначными и сложными ситуациями.

Важно отметить, что во всех экспериментах значения метрики Answer Faithfulness
были не ниже 0.88,
что говорит об отсутствии галлюцинаций и использовании языковыми моделями приведенного контекста.

Несмотря на все достижения,
у проведенного исследования есть ряд ограничений:
\begin{itemize}
    \item Размер и однотипность тестовых данных.
          Датасет содержит всего 5 подкорпусов, которые вместе
          насчитывают 63 заметки и 246 ссылок (таблица~\ref{tab:dataset_samples}).
          Их содержимое в основном~--- статьи Википедии.
          Это сильно ограничивает значимость метрик
          и делает невозможным оценку масштабируемости предложенных подходов
          и эксперименты с гиперпараметрами GraphRAG.
    \item Ограниченность онтологии.
          Набор используемых данных содержит всего несколько схожих тематических групп.
          Исходя из этого, невозможно предположить поведение системы на данных из другой области.
    \item Зависимость от качества разметки.
          Валидационных датасет создавался одним человеком,
          что вносит долю субъективности в типы некоторых связей.
\end{itemize}
